% =====================================================================
%  Harald Høffding, On Religious Freedom and Some Psychological and Ethical Questions Connected with It
%  Translated from the Danish: „Om Religionsfrihed og nogle dermed forbundne psykologiske og etiske Spørgsmaal“,
%  Tilskueren 3 (1886), June--July, pp. 448--459.
%
%  TRANSLATION -- GENERATED; DO NOT EDIT BY HAND.  Built by ebook-method/tilskueren-1886-90/tr_parts.py build --write
%  from the parts in ebook-method/tilskueren-1886-90/tr/ (GLOSSARY.md, TRANSLATOR-BRIEF.md).  Made from the collated
%  transcription texts/hoeffding/religionsfrihed-1886/transcription.tex (2026-10-15); [N] in the margin = the page of
%  Tilskueren.  Terms settled with Hans: Religionsfrihed = religious freedom, Trang = need, Følelsesliv = emotional life.
% =====================================================================
\documentclass[12pt,a4paper,openany]{book}
\usepackage[utf8]{inputenc}
\usepackage[T1]{fontenc}
\usepackage{libertinus}
\usepackage{libertinust1math}
\usepackage[german,english]{babel}
\usepackage[protrusion=true,expansion=true]{microtype}
\usepackage{setspace}
\usepackage[margin=1.3in]{geometry}
\usepackage{fancyhdr}
\usepackage{hyperref}
\hypersetup{pdftitle={On Religious Freedom}, pdfauthor={Harald Høffding}, hidelinks}
\pagestyle{fancy}
\fancyhf{}
\fancyhead[LE,RO]{\thepage}
\fancyhead[C]{\textit{H. Høffding: On Religious Freedom (1886)}}
\fancypagestyle{plain}{\fancyhf{}\fancyhead[LE,RO]{\thepage}}
\setstretch{1.3}
\usepackage{marginnote}
\usepackage{graphicx}
\renewcommand{\thefootnote}{\arabic{footnote})}
\newcommand{\opage}[1]{\marginnote{\footnotesize\textit{[#1]}}}
\newcommand{\sectionnum}[1]{\par\medskip\begin{center}#1\end{center}\smallskip}
\newcommand{\asterism}{\par\begin{center}*\quad*\\[-0.4em]*\end{center}}
\newcommand{\signature}[1]{\par\smallskip\hfill #1\hspace*{2em}\par}
\pdfstringdefDisableCommands{\renewcommand{\footnote}[2][]{}\renewcommand{\opage}[1]{}}

\begin{document}

\begin{titlepage}
\begin{center}
\vspace*{2em}
{\LARGE\bfseries \emph{TILSKUEREN.}}\\[1em]
{\large \emph{A MONTHLY}}\\[0.5em]
{\small FOR}\\[0.5em]
LITERATURE, SOCIAL QUESTIONS AND POPULAR\\
SCIENTIFIC ACCOUNTS.\\[1.5em]
{\small EDITED BY}\\[0.5em]
{\bfseries N. NEERGAARD.}\\[1.5em]
\emph{THIRD YEAR.}\\
\vfill
COPENHAGEN 1886.\\[0.4em]
{\bfseries P. G. PHILIPSEN.}\\[0.3em]
{\small THIELE'S PRINTING HOUSE.}\\[2em]
{\small Translated from the Danish}
\end{center}
\end{titlepage}
\clearpage

% --- p. 448 ---
\setcounter{footnote}{0}%
\opage{448}
\begin{center}{\large\bfseries On Religious Freedom and Some Psychological\\
and Ethical Questions Connected with It\footnote[1]{Essentially a lecture given in the Students' Association on 20 March 1886.}.}\\[0.5em]\rule{4em}{0.4pt}\end{center}\medskip

For the most part religious freedom is understood as something merely
\emph{external} and merely negative, namely what our Constitution expresses
by saying that no one may be deprived, on account of his confession of faith, of access to the full enjoyment
of civil or political rights. Even this is something very significant, when
one remembers what resistance the principle thus pronounced had
to struggle against before it prevailed. But for him who
reflects on the change that has been inaugurated, with the recognition of this
principle, in the whole spirit and direction of human life,
--- for him who holds that the outer and the inner may well come into
a certain opposition to one another, yet in the long run must correspond
to one another if life is to be healthy and vigorous, --- for him
who looks further than to the merely outward and provisional, --- for him
there lie in religious freedom psychological and ethical problems of
the greatest importance, problems which no one who takes life
seriously can avoid being confronted with. I shall here attempt
to bring out and illuminate some of these problems. In order that they may
properly emerge, I must begin with some historical, or, as
it is now the custom to say, sociological remarks.

\sectionnum{I.}

Positive religion is \emph{social} religion, in that it is built on
tradition, on the proclamation of the Word. In different forms
this principal feature is found again in all religions. It is not the
% --- p. 449 ---
\setcounter{footnote}{0}%
\opage{449}separate individuals who devise their religion, but they are initiated into
the one that lives on in the line of descent. As soon as one ascribes to the single individual
--- his need, his feeling, his power of judgment --- an absolute
right of decision in religious matters, one has already entered upon rationalism.
Religion is a \emph{historical} power; it is the proclamation
of a sacred history that forms the foundation on
which it is built.

In its simplest form religion is built on a tradition
that is propagated within one and the same family from generation to generation.
The spirits of the departed are worshipped with sacrifices and gifts, the hearth-fire is kept
burning in their honor, and their commands are respected as the highest law. In
the practice of this \emph{family religion} only those take part who belong
to the line of descent; the father of the family, who leads it, stands as the highest
living authority in religious as in other respects. To this
very day this kind of worship is the most widespread on
earth, in that it subsists beneath, alongside, or as a part of the
national religions. Only in Christianity and in Mohammedanism
has this worship disappeared entirely\footnote[1]{\emph{Henry Maine}: Early Law and Custom. p. 57 ff. --- It was \emph{Fustel de Coulange} who, in his interesting book ``La cité antique'', first demonstrated the great importance and primitive character of family religion within Greek and Roman religion.}.

\emph{The national religion} arose when the whole State received
its common gods and heroes, its common hearth.
Like the family, the State was originally both a religious and a political society.
There was no difference between State and Church, any more than between religious,
legal, and moral laws. It was the duty of every citizen
to take part in the public worship, but foreigners were excluded
from it. Each State had its gods, just as each family
had its own.

A significant turning point occurred when the \emph{idea of
the Church as a universal religious society} arose. Family differences
and national barriers were now to have only subordinate
validity. There shall no longer be Jew or Gentile,
Greek or barbarian. A purely spiritual society shall exist.
But here too the religious society, the spiritual unity, is conditioned
by the common transmission of a sacred history, from which
individuals draw the highest law for their lives and the deepest
assurance for their fate in this world and in a world beyond.
% --- p. 450 ---
\setcounter{footnote}{0}%
\opage{450}For European humanity, Christianity asserted
this idea in a victorious manner.

In the beginning the Church placed itself in sharp distinction from
other human societies as a kingdom that was not of this
world. It lived in the expectation that this world would
come to an end. When this expectation was not fulfilled, the Church began
its cultural-historical mission. It soon stamped its character upon inner and
outer life alike, whether its dominion appeared in
its ruling the State outright, or in the State's
seeing in it an indispensable instrument for securing peace and obedience.
After the medieval church-state, in which the Pope was the head of the Emperor,
came the rule of the ``territorial'' principle, by which the
religion of the sovereign was also to be that of his subjects (cujus
regio, ejus religio).

Then, from the eighteenth century onward, the principle of
the religious neutrality of the State emerges. What isolated thinkers
had demanded was now proclaimed as law for great nations.
The Constitution of the North American free states is the first that is built
on the principle of religious freedom.

A great good has thereby been attained. The spiritual struggle
can now be conducted more freely, since the external conditions are more equal for all.
However powerful prejudice and fanaticism may still continually be,
they no longer command such means as the stake and
fire or outward persecution.

But a great ideal, a great thought has been given up: the thought of
a common faith, of an inward spiritual fellowship for the whole
human race. And a great force seems bound to be lost, the one
that lies in union through a common view of life's
greatest problems. We now have leave each to be saved after his own
fashion. But will it not easily bring with it that we each sail our own sea
and lose mutual understanding?

\sectionnum{II.}

It will be answered from the side of the Church that there is no danger: the dissolution
will be only transitory, and on the ground of freedom
a new unity of faith, truer than the earlier enforced unity, will be attained.
And in the camp of the freethinkers many hope that the time is not
far off when all rational, honest, and cultivated people will have
turned their backs on the Church, so that at most only blockheads and
hypocrites will remain.

% --- p. 451 ---
\setcounter{footnote}{0}%
\opage{451}I share neither of these expectations. --- Just as little as
the Catholic Church, in spite of its great hope and its eager exertions,
will be able to win the Protestants back, just as little will
the Church be able to win the freethinkers back. Just as little as there is
in Catholicism a possibility of such a revival or such an
unfolding of spiritual superiority that the Protestants thereby
could be moved to give up their freedom of faith, their right to independent
religious examination, just as little is there any possibility
that the Church should remove from the world the motives which make people freethinkers.
The medieval Church was still able to
gather into itself all the spiritual strength of the age. But life is now too rich,
too independent, too manifold, for anything of the kind to be possible.
The great problems have been posed so clearly and sharply that
no dogmatic formula can cover them over. The need ---
whether one thinks that a solution of life's riddles is to be had
or not --- nonetheless to deal with them in one's own way has become
so powerful that it will not again be able to be subdued. It can be lulled
to sleep for a time, but will always awaken anew, as long as there exists
a higher spiritual life. And what holds for the sphere of knowledge
holds also for the life of imagination, of feeling, and of
will. Here too the free movement of the powers has already
become so great a good that one will not give it up again.
Life is lived no less deeply and richly outside the Church than in it.
Over against ecclesiastical faith there stands, and for the future will always
stand, a serious, free, and self-conscious view of life.

The freethinkers who think that the age of the positive religions
is past do not see that the Church will precisely gain by religious freedom.
The conflict that has been raised against the Church in the last centuries
concerns for the most part far more the Church's worldly power and
outward dominion than the dogmas. The conflict has not been so
much over the truth of the dogmas as over their right to determine
the civil and moral position of individuals. The experiences that
have been had in North America, in France after the Revolution, and in
Germany after 1875 prove that ecclesiastical interest and earnestness
increase with separation from the State. What apparently
weakens the Church in reality strengthens it. The free forces
are awakened to activity. Even the English state church, the stiffest
of all Protestant state churches, is built more and more on the principle
of voluntariness. Tocqueville derived the great power of religion in
North America from its having withdrawn entirely from politics.
% --- p. 452 ---
\setcounter{footnote}{0}%
\opage{452}There, in the promised land of freedom, it has so great an influence that
public opinion sets far narrower limits to free
inquiry and free views than it does here in Europe. Freethinkers
who are at the same time Democrats, and who work for the separation of the Church
from the State, must therefore have a great and cheerful faith in the
power of truth. --- What indeed would life be without such a faith?
But one can have it and yet hold that the end is very far off.

Since it is of importance for a question to which I shall
later draw attention, I shall here set forth, from a purely
psychological standpoint, some of the causes which in
my view will continue to bring adherence to positive
religious faith. There are of course very many different psychological
circumstances that operate here\footnote[1]{In my treatise ``Forholdet mellem Tro og Viden i dets historiske Udvikling'' (Vor Ungdom. 1885) I have had occasion to mention some causes which made themselves felt in the religious reaction at the beginning of this century (See especially p. 189---196).}; but there are especially two points
on which my experience and my investigations have led me to
lay stress.

The further spiritual development advances, the more
differentiated become the particular sides or directions of
spiritual life. Theoretical inquiry and practical interest, understanding
and feeling, science and art, separate themselves more or less
from one another. Each side of life poses its tasks in its own
way and solves them by its peculiar methods. The result
easily becomes a division of strength and interest, a lack of concentration,
which is unendurable to many natures. There will arise
a natural need to gather what is scattered into one point, or at least
to have one region where the distinction does not intervene, one sphere
where one can feel undivided life in its whole fullness. There are now
surely several different ways in which this need can be
satisfied. But this much is certain, that in the great folk religions
there is very much that meets it. These religions
have --- each in its time --- afforded full satisfaction for all
directions of spiritual life. They do not distinguish between
the true, the beautiful, and the good. They do not address themselves especially to
the understanding or to the imagination or to feeling or to the will,
but speak to all faculties at once and are understood only when they
% --- p. 453 ---
\setcounter{footnote}{0}%
\opage{453}can at once get them all to speak. Positive religion in
its classical form knows no special science, no
special art, no special ethics or politics. It always marks
a beginning rupture when one of these endeavors separates itself
as something apart and wants to go its own ways. For when, e.g.,
science becomes a province of its own beside religion, then
religion too becomes a province of its own beside science,
and is no longer the form and bearer of life in its totality.
The great positive religions are peculiar condensations or
syntheses of thoughts and ideals, hopes and endeavors,
which critical inquiry assigns to different provinces and
under different points of view. There will now be many whose
nature is so constituted that their most deeply reaching intellectual,
aesthetic, and ethical need demands satisfaction together, --- for
whom a thought that does not at the same time satisfy the imagination's need for
vivid images and feeling's need for hope and consolation is
no \emph{true} thought either. They cannot and will not distinguish (as we in
psychology distinguish) between feeling itself and the more or less
symbolic forms in which it comes to light,
or between a thought's ideal value and real validity. A religious
reaction often comes after a time poor in feeling and imagination, and
can then serve to enrich and strengthen the life of the spirit. The spiritual
life has, like all life, a rhythmic character, moves forward
through changing periods of gathering and scattering, unity and distinction,
synthesis and analysis. And there will always be natures who
are especially disposed for the one of these directions.

The second point that I wish to emphasize stands apparently
in conflict with the first, but only apparently. It is the
way in which life's greatest and strongest oppositions
come to outbreak and expression in the positive religions. Religious
feeling contains a need to experience exhaustively the
greatest suffering and the greatest delight, the greatest fall and the greatest
uprising that can be conceived. Only then does it think it knows life,
when it knows the greatest oppositions that life can bring; and
only when it knows life in this way does it trust that it can
overcome the world. In the great folk religions the most conflicting feelings
are awakened. Human minds are led through them,
as Dante was led through Hell, Purgatory, and Paradise.
Religious worship works especially through contrasts. In the
Eleusinian mysteries the participants seem to have celebrated the story of Demeter and
% --- p. 454 ---
\setcounter{footnote}{0}%
\opage{454}Persephone: the mother's grief at the daughter's disappearance,
her anguish of soul and her restless searching, finally her jubilation
at the rediscovery. So the Greeks too had their mater
dolorosa. In the year of the Christian Church the life of Christ is lived through anew
by the faithful, and with it a series of the greatest oppositions
in the domain of the emotional life. --- There are natures who find full nourishment for their need to
live through the oppositions only within
a positive religion. There is something in this need that
is common to all humanity, and that shows itself in every
serious person. The aesthetic effect of the tragic is connected
with this without, however, entirely coinciding with what
is here under discussion. Natures in whom this need asserts
itself with special strength will not easily distinguish, as others do,
between feeling itself and the symbols in which it gives itself expression
and finds rest; the symbol becomes for them dogma, the
perfect expression of the highest truth.

I mention this also in order to point out how little
mere logical and historical criticism in the religious domain
penetrates to the heart of the matter. Behind the logical contradictions and
historical improbabilities that criticism is able to point out
lie the mighty movements of feeling and oppositions of feeling.
It is these that make existence so great a riddle; how can
it contain the oppositions that call forth these conflicting
feelings? how can they spring from one and the same source?
how can they be reconciled and overcome? --- It is clear that one
may very well acknowledge and feel the reality of the oppositions
without approving the attempts to solve that riddle. One can recognize
that existence shows us both the greatest suffering and the
greatest happiness, both the deepest abasement and the most exalted
greatness, and yet not understand how they can be due to
the same cause, especially when this is supposed to be an absolutely good being.
Dogma stands here --- in my conviction --- just as powerless
as science. But the fully justified criticism of the solutions
ought not to bring about a failure to appreciate the strong impulses that lead to
the posing of the questions. For thereby one makes life narrow and poor.
We stand still --- despite all our philosophy of religion and our comparative
science of religion --- very far behind in regard
to understanding of the psychological foundation of the positive
religions.

% --- p. 455 ---
\setcounter{footnote}{0}%
\opage{455}I have strayed somewhat from my main subject and now return
to it.

\sectionnum{III.}

This much is certain, that the struggle will be a long one. In
the opinion of many it will become the eternal conflict. Should that be so, we must
make ourselves familiar with it. For a serious view of life
the whole of life is a great struggle, and this one will then be only one of its
many forms. But it is not a matter of indifference on what basis
and in what manner the conflict is waged. Is it to be waged from the ecclesiastical
side with the passage of Scripture as its motto with which an American
pastor took the field against Darwin and Spencer: ``Should I not hate
those who hate thee, Lord!''? And is it to be waged from the opposite side
with scorn and mockery? --- Against the worst excesses in this
direction provision has surely been made by the official recognition of religious freedom.
Outward, civil life can go its way, undisturbed
by the different faiths of the several individuals. The fanatical believers and
non-believers must in any case tolerate each other's presence in
society. But that is, as already remarked earlier, a decidedly
external and negative relation.

Nor will it be enough to point to the ever more
widespread and ever better grounded conviction that the religious
problems concern something which by its nature lies outside
the reach of strict science. For the question is precisely:
what then? Are we to let one another each sail his own sea, because
one can after all know nothing certain about these things? Will it not lead
to indifference or to a severing of every more deeply penetrating
spiritual connection between those who do not belong to the same sect?
--- And in any case one would then not be able to demand tolerance
of those who declare themselves to be altogether certain in their cause, ---
more certain than any science whatever could make them. The Catholics
thus often claim\footnote[1]{See e.g. a statement by a Catholic bishop, which is quoted in F. C. Baur: Kirchengeschichte des 19ten Jahrhunderts. P. 264.} that the Protestants \emph{must be} tolerant,
because they dare not maintain that their Church is the one true one,
but that the Catholics \emph{cannot be} so, because it is part of
their faith that their Church is the one true one. The Protestants perhaps
let the claim pass on to the freethinkers --- but the freethinkers
who think that all believers must be either stupid or
% --- p. 456 ---
\setcounter{footnote}{0}%
\opage{456}hypocritical, since they can believe what they believe, will send it back
with protest.

If religious freedom is to have more than external and negative significance,
we must investigate whether there may not be common domains
which can be withdrawn from the conflict, although they are of essential
influence on human life, and whether the possibility of mutual
understanding and sympathy is altogether cut off in respect of the very questions
about which the conflict is being waged. In the first respect
it would be of great importance above all if there is a humane ethics
building on a natural human basis. In the latter respect
it will depend on whether people who stand on different
religious standpoints can nevertheless have a certain spiritual
kinship with one another.

If the only point of view which the different religious standpoints
know in relation to one another is: truth or lie, then
it is of no use to investigate these questions. But in that case
it is also hard to see how one can accept
external, civil religious freedom. How can one put up with it that
liars exercise civil rights and take part in deciding
public affairs? And it is also hard to see how
one can recognize the independence of science. If there is a
science which reaches its goal without regard to and without guidance
from religious views, then at least a part of the truth
can after all be attained by human means. For however one may
judge between knowledge and faith, --- in the end there can be only
one truth. --- But should the capacity of human nature to
develop by its own powers be confined to civil
life and to science? Is there not also a purely
human ethics which, at any rate to a certain degree, is independent
of religious assumptions? --- I do not understand how
anyone who sees in every human being a distinctive personal
being can deny this. A personal being can only have in itself
the ultimate standard for what it is to recognize as true,
noble and good. Whether it be something supernatural or something
natural that personality stands over against: in its own inner self it must
seek the principles for understanding and appraisal of everything that
presents itself to it. Whatever ideals man may then recognize,
and however exalted the ethical law may stand even \emph{for him}, they must
after all be ideal and law for him, and become so only by
finding points of attachment in his inner spiritual nature. It is
% --- p. 457 ---
\setcounter{footnote}{0}%
\opage{457}the great significance of Socrates's \emph{Command: Know thyself}! by
which he became the founder of ethics. Philosophical ethics in
its various forms has built on this principle, which may be called
\emph{the principle of personality}, the principle both of free inquiry
and of free conscience. In working with
this principle, philosophical ethics at the same time draws the consequences of the
principle of universal love of humankind. It cannot accept
the view that the difference between acceptance or
non-acceptance of certain religious doctrines is to draw an absolute
boundary between human beings, to divide them eternally into two parts. As love of humankind
has broken through the barriers that separated nation
from nation and class from class, so it will and must also break through
the barriers which confessional differences have set up
against the full recognition of the personal life in all
human beings. The religious conflict will not cease thereby,
and it will not lose its sharpness. Nothing is to be abated
in the question of what truth is. But the conflict can be
sharp without being bitter and hateful.

If believers can regard the non-believers'
virtues only as splendid vices, --- if the non-acceptance of the Church's
doctrine is to be derived solely from evil will and antipathy to holiness,
--- if the conflict between faith and unbelief is wholly one with
the conflict between good and evil, then all hope is gone. Often it sounds
as if that were so. But much of what is said is also
only big words, not thoughts, at any rate not thoughts whose consequences
one draws. For the full consequence of such
thoughts would lead far back --- back to the Middle Ages, to the
time that lies before the recognition of religious freedom and
the development of modern culture!

For the believer the purely human ethics will of course
not be able to be exhaustive. But it must nevertheless be of interest
also to him to see how far one can reach along this path,
and he will quite well be able himself to take part in the inquiry.
He will reserve the right to set his own spire on the building.
But there will in any case be a protected domain of great importance,
into which the religious struggle need not penetrate, just as
little as it now penetrates the territories of physics and chemistry\footnote[1]{Concerning the nature and significance of humane ethics I permit myself to refer partly to my work ``Om Grundlaget for den humane Etik'' (\opage{458}Copenhagen 1876), partly to my treatise ``Om den filosofiske Etiks Principer'' in ``Oversigt over det kgl. danske Videnskabernes Selskabs Forhandlinger i Aaret 1886''.}.

% --- p. 458 ---
\setcounter{footnote}{0}%
\opage{458}Only when it is recognized that ethics has its basis in
human nature itself will actual, positive religious freedom
become possible, since the greatest obstacle will have been cleared away for
mutual understanding and sympathy between people who stand at
different religious standpoints.

There is yet another important condition which must be present,
and which --- partly on account of the dogmatic manner in which
the religious struggle, both as attack and as defense,
has been waged --- is as yet developed only to a slight degree among us.
It is the sense for personal life, especially the emotional life, in
all its different forms. It is also very closely connected
with the faith that in every human being there stirs, or at
any rate there is a disposition to, a distinctive life. It presupposes
a delight in hearing each bird sing with its own beak. But
it demands a special education. We are as yet far too dogmatic
in our psychology. We proceed far too confidently from the assumption that others
perceive, think and feel in quite the same way as we. Individual
differences make themselves felt even in the simplest
mental activities, differences which one has only just
begun to bring to light. How much more must
the higher spiritual life take shape differently, and how careful
we must be not to draw hasty conclusions from the
manner in which the emotional life gives expression in words or symbols!
These may sound alike and yet cover inner movements
of very different kinds. Only he who has the capacity to distinguish
between shell and kernel discovers these differences, this
richness. He who stops at the positive or negative confession of faith
gets no further than the surface. The flower
that grows in the tall grass is discovered only by him who bends the
long blades aside. But he who does so will not merely find
great differences and rich nuances where he did not suspect them.
He will also discover that spiritual kinship is other
than dogmatic kinship. Just as behind the same confession of faith
a most different inner emotional life may lie, so behind
different confessions of faith a closely akin emotional life may lie.
It may be difficult to penetrate to it, and we stand far
% --- p. 459 ---
\setcounter{footnote}{0}%
\opage{459}behind in this art. But until one learns it, one has not yet
learned to know what is essential in man. As Schiller has already said:
what we think belongs to all; only what we feel
is our own. And the emotional life always has a certain independence of
thoughts and ideas. Let us take a single example. For
the serious believer as for the serious freethinker life stands
as a struggle in which it is a matter of upholding the ideal; they
feel the tragic in this struggle and the great antitheses through
which it leads; their feeling of life, their feeling about existence
is in the end determined by these experiences. Should there not be
greater spiritual kinship between two such persons than
between such as belong to the same confession but whose emotional life
moves in wholly different courses? --- Psychologically viewed, the difference
between the believer and the freethinker is often only this, that the former
has his dogmatic forms in which he finds expression for his
innermost experiences of life, whereas the freethinker must say with Faust:
\begin{verse}\small
Ich habe keinen Namen dafür!\\
Gefühl ich alles,\\
Name ist Schall und Rauch.
\end{verse}

Amid the standing, perhaps eternal, conflict it would be a
great good if an understanding of the kind I have here indicated
could develop. It would be not only an enrichment and extension
of our life, it would be the only possible form of
the spiritual unity of the human race, now that the great idea of
uniting all under one confession has had to succumb. Religious freedom
would then be not merely something external and negative,
but without abating anything in what the great conflict essentially
turns on, the gathering and unifying forces would be able to work.
\signature{\textsc{Harald Høffding}.}

\end{document}
